\subsection{不可约元素}

定义 6.——序群 G 的元素 x 称为不可约的，如果它是 G 的严格正元素集合中的极小元素。

设 x 是序群 G 的一个不可约元；若 y 是 G 的正元素，则元素 \(\inf(x, y)\) 若存在，只能等于 x 或 0。因此在格序群 G 中，每个正元素 y 或者大于不可约元 x，或者与 x 互素；特别地，两个不同的不可约元是互素的。

\begin{enumerate}
\def\labelenumi{(\Roman{enumi})}
\setcounter{enumi}{503}
\tightlist
\item
  元素 \(p\)（属于 \(\mathbf{A}\)）称为不可约的，如果理想 \((p)\) 是序群 \(\mathcal{P}^*\) 的一个不可约元；这表示 \(p\) 既非零亦非可逆元，且 \(\mathbf{A}\) 中任意整除 \(p\) 的元素与 \(p\) 或与 1 相伴。如果 \(\mathcal{P}^*\) 是格序的，那么每个 \(a \in \mathbf{A}\) 要么与 \(p\) 互素，要么是 \(p\) 的倍元。
\end{enumerate}

例子（DIV）。—— 1）整数 \(p > 0\) 在 \(\mathbf{Z}\) 中不可约，当且仅当它是素数（I，第 50 页）。

\begin{enumerate}
\def\labelenumi{\arabic{enumi})}
\setcounter{enumi}{1}
\tightlist
\item
  域 K 上关于单个未定元的多项式是不可约的，当且仅当它在通常意义下不可约（IV，第 13 页）。
\end{enumerate}

命题 14.——要使 \(x > 0\) （序群 \(G\) 的元素）成为不可约元，只需它满足以下性质：

\begin{enumerate}
\def\labelenumi{(\Alph{enumi})}
\setcounter{enumi}{15}
\tightlist
\item
  关系 \(x \leqslant y + z\) , \(y \geqslant 0\) , \(z \geqslant 0\) 蕴含 \(x \leqslant y\) 或 \(x \leqslant z\) 。
\end{enumerate}

当 G 是格序群时，这个条件是必要的。

如果 G 是格序的且 x 不可约，我们刚刚看到，y 要么大于 x，要么与 x 互素；在后一情形，VI 第 15 页的推论 2 表明 z 大于 x。反之，设该条件成立：若 \(0 \leqslant y \leqslant x\) ，则令 \(x = y + z\) ( \(z \geqslant 0\) )，可得 \(x \leqslant y\) 或 \(x \leqslant z\) ；在第一种情形有 x = y；在第二种情形有 \(x \leqslant x - y\) ，故 \(y \leqslant 0\) ，从而 y = 0；这表明 x 确实是不可约的。

命题 14（DIV）。——要使 A 中非零元素 p 不可约，只需：p 不是单位，且 p 不会在整除 A 中两个元素之积而不整除其中至少一个。当 \(P^{*}\) 是格序的时，此条件也是必要的。

注记。——命题 14（DIV）也可表述如下：若 p 是 A 中非零元素，使理想 \((p)\) 是素理想（I，第 117 页，定义 3），则 p 不可约；反之，若 \(P^{*}\) 是格序的且 p 不可约，则理想 \((p)\) 是素理想。

命题 15.——在序群 G 中，设 \((p_{\iota})_{\iota \in I}\) 是 G 中满足条件 \((P)\) （从而不可约）的两两不同的正元素族。则映射

\[
(n _ {i}) _ {i \in I} \mapsto \sum_ {i \in I} n _ {i} p _ {i}
\]

是序群 \(\mathbf{Z}^{(I)}\) （序群 Z 的直和，VI，第 7 页）到由 \(p_{i}\)

生成的 G 的序子群上的一个同构。只需证明关系式 \(\sum_{\iota \in I} n_{\iota} p_{\iota} \geqslant 0\) 等价于 \(n_{\iota} \geqslant 0\) 对所有 \(\iota\) 成立，

特别地，关系式 \(\sum_{\iota \in I} n_{\iota} p_{\iota} = 0\) 将蕴含 \(n_{\iota} = 0\) 对所有 \(\iota\) 成立，因此这将

表明族 \((p_{\iota})\) 线性无关。现设 \(I'\) （相应地 \(I''\) ）为 \(I\) 的有限子集，由那些 \(\iota\) （使 \(n_{\iota} > 0\) ，相应地使 \(n_{\iota} < 0\) ）组成；我们有

\[
\sum_ {i \in I ^ {\prime}} n _ {i} p _ {i} \geqslant \sum_ {i \in I ^ {\prime \prime}} (- n _ {i}) p _ {i}.
\]

特别地，对 \(\lambda \in I''\) ，这蕴含 \(p_{\lambda} \leqslant \sum_{i \in I'} n_i p_i\) ，并且由归纳法

从性质 (P) 可知，必有 \(p_{\lambda} \leqslant p_{\iota}\) 对某个 \(\iota \in I'\) 成立；由于 \(p_{\iota}\) 不可约，这将蕴含 \(p_{\lambda} = p_{\iota}\) ，这是荒谬的。因此 \(I''\) 为空，这就证明了命题。

定理 2. --- 设 G 是一个滤过群。则以下性质等价：

\begin{enumerate}
\def\labelenumi{\alph{enumi})}
\item
  G 同构于形如 \(\mathbf{Z}^{(I)}\) 的序群。
\item
  G 是格序的，并满足以下条件：
\end{enumerate}

（MIN）\(\mathbf{G}\) 中由正元素组成的每个非空集合都有一个极小元素。

\begin{enumerate}
\def\labelenumi{\alph{enumi})}
\setcounter{enumi}{2}
\item
  G 满足条件 (MIN)，且 G 的每个不可约元都具有性质 (P)。
\item
  G 由其不可约元生成，且 G 的每个不可约元都具有性质 (P)。
\end{enumerate}

首先证明 \(a\) 蕴含 \(b\) 。群 \(\mathbf{Z}^{(I)}\) 是格序的，因为它是全序群的直和。另一方面，设 E 为 \(\mathbf{Z}^{(I)}\) 中正元素组成的一个非空集合，且令 \(x = \sum_{i} n_i e_i\) 为 E 的一个元素（其中

\((e_{i})\) 表示 \(\mathbf{Z}^{(I)}\) 的自然基）；只有有限个（数目为 \(\prod_{i}(n_{i} + 1)\) ）

元素 \(y\)（在 \(\mathbf{Z}^{(I)}\) 中）满足 \(0 \leqslant y \leqslant x\) ，因此集合 \(F\) （由 \(E\) 中小于或等于 \(x\) 的元素组成）更是有限的；由于它非空，它包含一个极小元素（集合论，III，第 170 页，推论 2），这显然是 \(E\) 的一个极小元素。

显然 \(b)\) 蕴含 \(c)\) ，由命题 14。让我们证明 \(c)\) 蕴含 \(d)\) 。由于 \(G\) 是滤过的，只需（VI，第 4 页，命题 4）验证：等于不可约元之和的正元素所成的集合 \(F\) （\(G\) 的正元素之集的子集）等于 \(G_{+} - \{0\}\) 。若非如此，则由 (MIN)，\(F\) 在 \(G_{+} - \{0\}\) 中的补集将有一个极小元素 \(a\) ；按定义 \(a\) 不是不可约的，故而是两个严格正元素 \(x\) 与 \(y\) 之和；由于 \(x < a\) 且 \(y < a\) ，这些元素属于 \(F\) ，因而是不可约元之和，由此可知 \(a\) 也是，这与假设矛盾。最后，\(d)\) 蕴含 \(a)\) ，由命题 15。

我们将把定理 2 应用于主理想整环（VII，第 4 页）和唯一分解整环（AC，VII，第 3 节）中的整除性理论，以及戴德金环（AC，VII，第 2 节）中理想的研究。

